% Source: https://artofproblemsolving.com/community/c3h1664765_some_questions_from_turkey_maths_olympiads_2015
% Stable identity: aops-topic-1664765-q01
\textbf{At most $4$.} Note $2015=5\cdot13\cdot31$. Every other element shares a prime factor with $2015$, so it is divisible by $5$, $13$ or $31$. No prime can divide three elements, and each of $5,13,31$ already divides $2015$; so each of these primes divides at most one other element. Hence there are at most $3$ elements besides $2015$, and at most $4$ in all.

\textbf{$4$ is possible.} Take
\[
2015=5\cdot13\cdot31,\qquad a=2\cdot3\cdot5=30,\qquad b=2\cdot7\cdot13=182,\qquad c=3\cdot7\cdot31=651 .
\]
Each pair shares a prime ($2015$ shares $5$, $13$, $31$ with $a$, $b$, $c$; $a,b$ share $2$; $a,c$ share $3$; $b,c$ share $7$). Each of the primes $2,3,5,7,13,31$ divides exactly two of the four numbers, so no three have a common prime factor.

The answer is $\boxed{\textbf{(C)}\ 4}$.
